\documentclass[11pt]{article}
\usepackage{amsmath, amssymb, amsthm}
\usepackage[shortlabels]{enumitem}
\usepackage[margin=1in]{geometry}
\usepackage{xfrac}

%%
%% NOTE: This file uses macros that would normally be defined in an
%% external style file (e.g., worksheet-setup.tex or course-macros.sty).
%% These macros are NOT defined in this preamble, so Pandoc will not
%% recognize them. They include:
%%
%%   \R, \Z, \N, \C, \Q   → blackboard-bold number sets
%%   \ds                   → \displaystyle
%%   \Span                 → \operatorname{span}
%%   \rank                 → \operatorname{rank}
%%   \rref                 → \operatorname{rref}
%%   \im                   → \operatorname{im}
%%   \sgn                  → \operatorname{sgn}
%%   \grstep               → Gaussian row reduction arrow
%%   \amatrix              → augmented matrix environment
%%   \nicefrac             → \sfrac (from xfrac)
%%   \blankbox             → blank answer box
%%

\newtheorem*{theorem}{Theorem}
\newtheorem*{definition}{Definition}

\pagestyle{empty}

\begin{document}

\begin{center}
  {\large\textbf{Math 22a --- Linear Algebra}} \\[4pt]
  \textbf{Problem Set 8: Vector Spaces and Eigenvalues} \\[2pt]
  Fall 2025
\end{center}

\bigskip

\begin{enumerate}

\item Let $V = \R^3$ and let $W = \{(x, y, z) \in \R^3 : 2x - y + 3z = 0\}$.
\begin{enumerate}[(a)]
  \item Show that $W$ is a subspace of $\R^3$.
  \item Find a basis for $W$.
  \item What is $\dim(W)$?
\end{enumerate}

\item Consider the following vectors in $\R^4$:
\[
  \mathbf{v}_1 = \begin{pmatrix} 1 \\ 0 \\ 2 \\ 1 \end{pmatrix}, \quad
  \mathbf{v}_2 = \begin{pmatrix} 0 \\ 1 \\ -1 \\ 3 \end{pmatrix}, \quad
  \mathbf{v}_3 = \begin{pmatrix} 2 \\ 1 \\ 3 \\ 5 \end{pmatrix}, \quad
  \mathbf{v}_4 = \begin{pmatrix} 1 \\ 1 \\ 1 \\ 4 \end{pmatrix}.
\]
\begin{enumerate}[(a)]
  \item Find $\Span\{\mathbf{v}_1, \mathbf{v}_2, \mathbf{v}_3, \mathbf{v}_4\}$ by row reducing and identifying pivot columns.
  \item What is $\rank(A)$ where $A = [\mathbf{v}_1 \; \mathbf{v}_2 \; \mathbf{v}_3 \; \mathbf{v}_4]$?
  \item Find a basis for $\im(A)$ and a basis for $\ker(A)$.
\end{enumerate}

\item Row reduce the following augmented matrix to $\rref$:
\[
  \begin{amatrix}{3}
    1 & 2 & -1 & 4 \\
    2 & 5 & 0  & 9 \\
    -1 & 0 & 3 & 1
  \end{amatrix}
  \grstep{R_2 - 2R_1}
  \begin{amatrix}{3}
    1 & 2 & -1 & 4 \\
    0 & 1 & 2  & 1 \\
    -1 & 0 & 3 & 1
  \end{amatrix}
  \grstep{R_3 + R_1}
  \begin{amatrix}{3}
    1 & 2 & -1 & 4 \\
    0 & 1 & 2  & 1 \\
    0 & 2 & 2  & 5
  \end{amatrix}
\]
Continue to $\rref$ and solve the system $A\mathbf{x} = \mathbf{b}$.

\item Let $A = \begin{pmatrix} 3 & 1 \\ 0 & 2 \end{pmatrix}$.
\begin{enumerate}[(a)]
  \item Find the eigenvalues of $A$.
  \item For each eigenvalue $\lambda$, find the eigenspace $E_\lambda = \ker(A - \lambda I)$.
  \item Is $A$ diagonalizable? If so, find matrices $P$ and $D$ such that $A = PDP^{-1}$.
\end{enumerate}

\item Let $T \colon \R^3 \to \R^2$ be the linear transformation with standard matrix
\[
  A = \begin{pmatrix} 1 & -2 & 3 \\ 2 & 0 & -1 \end{pmatrix}.
\]
\begin{enumerate}[(a)]
  \item Find $\ker(T)$. What is $\dim(\ker(T))$?
  \item Find $\im(T)$. What is $\rank(T)$?
  \item Verify the Rank-Nullity Theorem: $\rank(T) + \dim(\ker(T)) = \dim(\R^3)$.
\end{enumerate}

\item Let $f \colon \Z \to \Z$ be defined by $f(n) = 3n + 1$. Determine whether $f$ is injective, surjective, and/or bijective. Note that in $\Q$ and $\R$ the story may differ --- explain why.

\item Show that if $\lambda$ is an eigenvalue of an invertible matrix $A$, then $\nicefrac{1}{\lambda}$ is an eigenvalue of $A^{-1}$.

\item Compute $\ds \sum_{k=0}^{n} \binom{n}{k} (-1)^k$ and prove your answer by induction. Also compute $\ds \sum_{k=1}^{n} k \cdot k!$ and find a closed form.

\item The \textbf{sign} of a permutation $\sigma$ is $\sgn(\sigma) = (-1)^{m}$ where $m$ is the number of inversions. Compute $\sgn(\sigma)$ for:
\begin{enumerate}[(a)]
  \item $\sigma = (1 \; 2 \; 3)$
  \item $\sigma = (1 \; 3)(2 \; 4)$
\end{enumerate}

\item Fill in the blank: \blankbox{2in}{0.5in}

A linear map $T \colon V \to W$ between finite-dimensional vector spaces over $\C$ satisfies the Rank-Nullity Theorem regardless of whether $V$ and $W$ are over $\R$, $\Q$, or $\C$.

\end{enumerate}

\end{document}
